\documentclass[a4paper]{article}
\usepackage{amsfonts}
\usepackage{amsmath}
\usepackage{amssymb}
\usepackage{graphicx}

\begin{document}

\noindent \large Auteur: Victoria Mazur \\
\noindent \large Studierichting: Informatica\\
\noindent \large Oefeningenreeks 4A nr. 20.7\\

\medskip

\normalsize

$\displaystyle\int \sinh^2 x dx$ \\

Stel $\left\{
\begin{array}{l}
u = \sinh x \\
dv = \sinh x dx
\end{array}
\right. \Rightarrow \left\{
\begin{array}{l}
du = \cosh x dx \\
v = \cosh x
\end{array}
\right.$ \\

Parti�le integratie: $\displaystyle\int u dv = uv - \displaystyle\int v du$ \\

$= \sinh x \cosh x - \displaystyle\int \cosh x \cdot \cosh x dx$ \\

$= \sinh x \cosh x - \displaystyle\int \cosh^2 x dx$ \\

We weten dat $\cosh^2 x - \sinh^2 x = 1$, dus $\cosh^2 x = 1 + \sinh^2 x$ \\

$= \sinh x \cosh x - \displaystyle\int \left(1 + \sinh^2 x\right) dx$ \\

$= \sinh x \cosh x - \displaystyle\int 1 dx - \displaystyle\int \sinh^2 x dx$ \\

$= \sinh x \cosh x - x - \displaystyle\int \sinh^2 x dx$ \\

$\displaystyle\int \sinh^2 x dx = \sinh x \cosh x - x - \displaystyle\int \sinh^2 x dx$ \\

$\displaystyle\int \sinh^2 x dx + \displaystyle\int \sinh^2 x dx = \sinh x \cosh x - x$ \\

$2 \displaystyle\int \sinh^2 x dx = \sinh x \cosh x - x + C$ \\

$\displaystyle\int \sinh^2 x dx = \dfrac{\sinh x \cosh x - x}{2} + C$ \\

$\displaystyle\int \sinh^2 x dx = \dfrac{\sinh x \cosh x}{2} - \dfrac{x}{2} + C$ \\

\begin{thebibliography}{999}
%\bibitem{a} Referentie
\end{thebibliography}
\end{document}